\documentclass[11pt]{article}

% Self-contained and pdflatex-safe: no external style files, no native Bengali
% glyphs in the source, so it compiles on arXiv out of the box. For a venue
% submission, reformat with the official ACL or LREC style. Bengali examples are
% given in transliteration to keep the build font-independent.

\usepackage[T1]{fontenc}
\usepackage[utf8]{inputenc}
\usepackage[margin=1in]{geometry}
\usepackage{booktabs}
\usepackage{amsmath}
\usepackage{microtype}
\usepackage{times}
\usepackage[hidelinks]{hyperref}
\usepackage{url}

\title{\textbf{A Bengali-First, Grapheme-Cluster-Aware Tokenizer with Zero Conjunct Fragmentation}}
\author{
  Konko Maji \\
  Project Bornomala, West Bengal, India \\
  \texttt{work.konkomaji@gmail.com} \\
  ORCID: 0009-0008-3561-3612
}
\date{2026}

\begin{document}
\maketitle

\begin{abstract}
Every large language model reads text through a tokenizer, and nearly every
mainstream tokenizer is optimised for English. For Bengali, an abugida whose
written unit (the grapheme cluster) spans several codepoints but reads as one
symbol, character-level subword methods routinely place token boundaries inside
a cluster, splitting conjuncts into fragments that correspond to nothing a reader
recognises. We present a Bengali-first tokenizer that makes this impossible by
construction. The method normalises to NFC, segments into Unicode UAX~\#29
grapheme clusters, remaps each cluster to an atomic symbol, and trains a subword
model over atoms, so every learned token is a whole number of clusters. On a
held-out Bengali evaluation set, our tokenizer attains a token fertility of
$1.39$ tokens per word, ahead of AI4Bharat IndicBERTv2 ($1.52$), Sarvam-1
($2.01$), GPT-4o ($2.61$), multilingual BERT ($3.01$), and DeepSeek-V3 ($3.02$),
while fragmenting $0.06\%$ of grapheme clusters against $3.6\%$ to $29.9\%$ for
the baselines. Training is CPU-only. The tokenizer, the evaluation code, and the
raw numbers are released under Apache~2.0, and the comparison reproduces with a
single command.
\end{abstract}

\section{Introduction}
A widespread assumption holds that a multilingual model treats its supported
languages roughly equally. It does not. Before a single transformer layer runs,
text passes through a tokenizer whose vocabulary was induced from a mostly
English mixture. A language whose characters and scripts are rare in that mixture
is split into more, smaller pieces. This is measured by \emph{fertility}, the
average number of subword tokens per word, and it has two consequences: a direct
cost, since more tokens mean more money and less usable context, and a capability
cost, since a fixed context window then holds less content.

For Bengali the problem is not only quantitative. Bengali is an abugida: a
written unit is a base consonant carrying an inherent vowel, modified by vowel
signs and combined with other consonants through the virama into conjunct
ligatures. Such a unit, the grapheme cluster, spans one to roughly eight
codepoints, yet a reader perceives it as a single symbol. Character-level Byte
Pair Encoding (BPE) and Unigram, the default methods, operate on codepoints and
can therefore place a token boundary in the middle of a conjunct. The resulting
token corresponds to no readable unit, and error metrics computed on codepoints
understate the damage.

We present a tokenizer designed so that a token can never split a grapheme
cluster. Our contributions are: (1) a grapheme-atom scheme that makes conjunct
fragmentation structurally impossible for covered clusters; (2) the first
measured, reproducible comparison of how efficiently mainstream tokenizers encode
Bengali, on identical held-out text; and (3) an open, CPU-only, Apache-2.0
release of the tokenizer and the evaluation.

\section{Background and Related Work}
\paragraph{Tokenization and fertility.} Fertility is the ratio of subword tokens
to whitespace words; lower is better. Purpose-built Indic tokenizers reduce it
substantially relative to English-first models: Sarvam-1 reports fertility in the
range $1.4$ to $2.1$ across ten Indic languages, and IndicSuperTokenizer improves
average fertility over strong baselines through linguistically grounded
pre-tokenization and a subword-to-superword stage. A complementary metric, the
single-token retention rate (STRR), reports the fraction of words preserved as
exactly one token and exposes vocabulary allocation that fertility alone hides.

\paragraph{Bengali script and grapheme clusters.} The correct unit for Bengali is
the extended grapheme cluster of Unicode UAX~\#29, not the codepoint. Prior work
on Bengali handwriting recognition (the BnGraphemizer line) found that
character-based tokenizers fail because conjunct glyphs have visual forms
unrelated to their constituent characters, and moved to grapheme-level
tokenization. We adopt the same principle for text tokenization and make the
grapheme cluster the atomic unit of the subword model itself. All processing
follows NFC normalisation (UAX~\#15), without which the same visible word admits
several codepoint sequences.

\paragraph{The gap.} No published tokenizer has been induced from Bengali alone
with a grapheme-cluster-respecting design, and, to our knowledge, Bengali has not
appeared in a single published cross-tokenizer fertility comparison. We address
both.

\section{Method}
The tokenizer is a pipeline of five stages.

\paragraph{Normalisation.} Input is put in NFC and passed through a documented
zero-width joiner and non-joiner policy that preserves those characters by
default, since they carry ligature-formation intent.

\paragraph{Grapheme segmentation.} Text is segmented into UAX~\#29 extended
grapheme clusters. This is the unit a Bengali reader perceives as one symbol.

\paragraph{Atom remapping.} Each grapheme cluster is mapped to a single symbol
drawn from the Unicode Private Use Area, called an atom. Frequent clusters each
receive their own atom; every individual codepoint also receives one, so a rare
cluster decomposes into codepoint atoms rather than becoming unknown. The map is
reversible and is stored with the tokenizer.

\paragraph{Subword training.} A BPE or Unigram model is trained over atom
strings. Because every atom is a whole cluster, every merge combines whole
clusters, and every learned token is a sequence of whole clusters. Word
boundaries are carried by a Metaspace marker. Crucially, all atoms are forced
into the base vocabulary, which guarantees that any covered cluster is a valid
single token, so encode followed by decode reproduces the exact normalised text.

\paragraph{Coverage guarantee.} The whole Bengali Unicode block and ASCII are
guaranteed atoms, so any Bengali or code-mixed English text round-trips
regardless of the training corpus. Conjunct fragmentation is therefore $0$ for
every cluster that has an atom; only clusters below the atom-frequency threshold
can fragment, and their rate is measured and reported, not hidden.

\section{Experimental Setup}
\paragraph{Training data.} The tokenizer is trained on the first $12{,}000$
articles of the Bengali Wikipedia dump (\texttt{wikimedia/wikipedia},
configuration \texttt{20231101.bn}), a BPE vocabulary of $32{,}000$. Training is
CPU-only.

\paragraph{Evaluation data.} All tokenizers are evaluated on $878$ held-out lines
taken from Bengali Wikipedia articles \emph{after} the first $12{,}000$, so the
text was never seen during our training. Every tokenizer is measured on this same
text after NFC normalisation.

\paragraph{Baselines.} We compare against the real, official public tokenizers of
AI4Bharat IndicBERTv2, Sarvam-1, XLM-RoBERTa, multilingual BERT, DeepSeek-V3, and
GPT-4o (the OpenAI \texttt{o200k} encoding).

\paragraph{Metrics.} Fertility is tokens per whitespace word. STRR is the
fraction of words kept as one token. Bytes per token is UTF-8 bytes over tokens.
Conjunct fragmentation is the fraction of Bengali grapheme clusters that a token
boundary splits, computed from each tokenizer's own character offsets; GPT-4o
exposes no offsets, so its value is not measured rather than estimated.

\section{Results}
Table~\ref{tab:main} reports the comparison, sorted by fertility.

\begin{table}[t]
\centering
\small
\begin{tabular}{lrrrr}
\toprule
Tokenizer & Fert. & STRR & B/tok & Conj. frag. \\
\midrule
\textbf{Ours (BPE 32k)} & \textbf{1.390} & \textbf{0.766} & \textbf{13.03} & \textbf{0.0006} \\
IndicBERTv2 (AI4Bharat) & 1.520 & 0.669 & 11.92 & 0.0364 \\
Sarvam-1 (Sarvam AI) & 2.005 & 0.490 & 9.04 & 0.0942 \\
XLM-RoBERTa (Meta) & 2.351 & 0.406 & 7.71 & 0.1034 \\
GPT-4o (OpenAI) & 2.608 & 0.095 & 6.95 & n/a \\
mBERT (Google) & 3.012 & 0.317 & 6.02 & 0.2164 \\
DeepSeek-V3 & 3.024 & 0.071 & 5.99 & 0.2988 \\
\bottomrule
\end{tabular}
\caption{Cross-tokenizer comparison on held-out Bengali. Lower fertility and
lower fragmentation are better; higher STRR and bytes per token are better. Best
in bold.}
\label{tab:main}
\end{table}

Our tokenizer attains the lowest fertility ($1.39$), ahead of AI4Bharat's own
IndicBERTv2 ($1.52$) and well ahead of Sarvam-1 ($2.01$) and the general models.
It retains the most whole words (STRR $0.77$) and packs the most bytes per token
($13.03$). It does this with a $32{,}000$ vocabulary, smaller than several
baselines, so the efficiency is per vocabulary slot as well as absolute.

\section{Analysis: conjunct integrity}
The decisive difference is conjunct fragmentation. Our tokenizer splits
$0.06\%$ of grapheme clusters, all of them rare clusters below the atom-frequency
threshold. The baselines split between $3.6\%$ (IndicBERTv2) and $29.9\%$
(DeepSeek-V3). A split conjunct is not a benign subword boundary: it severs the
written unit a reader recognises, and it does so invisibly to codepoint-level
metrics. A Bengali-first design that treats the grapheme cluster as atomic
removes this failure mode by construction rather than reducing it by tuning.

\section{Limitations}
The checked-in artifact is trained on Wikipedia, which is not weighted toward the
literary and formal register where Bengali conjunct density is highest; the final
tokenizer trains on a literary-weighted corpus. Fragmentation is near zero but not
exactly zero, because clusters below the frequency threshold decompose; setting
the threshold to one drives it to zero for all seen clusters at the cost of a
larger atom set. The held-out set is Wikipedia and is unseen relative to training,
but is neither literary nor dialectal; a literary and a dialect evaluation set are
planned. Decode normalises surrounding whitespace, which subword tokenizers
generally do not carry. Vocabulary sizes differ across systems, which we note
rather than control, since a smaller vocabulary leading a larger one is part of
the result.

\section{Ethics and Broader Impact}
This work is part of Project Bornomala, a non-commercial research effort based in
West Bengal to build a Bengali-first, dialect-aware language model and to preserve
the Bengali language and its dialects. Reducing the token cost of Bengali is a
language-equity contribution: it lowers the price and raises the usable context
for a language spoken by more than 250 million people. All artifacts are open, and
every reported number states the data behind it; no estimate is presented as a
measurement.

\section{Reproducibility}
The tokenizer, the evaluation harness, the trained artifact, and the raw numbers
are released under Apache~2.0. The comparison reproduces with a single command,
\texttt{python scripts/compare.py --tokenizer artifacts/bn-bpe-32k --skip 12000
--limit 800}, on CPU.

\section{Conclusion}
Treating the grapheme cluster, not the codepoint, as the atomic unit of a subword
tokenizer yields a Bengali tokenizer that is both more efficient than the
incumbents and structurally incapable of breaking a conjunct. We hope the open
release and the first reproducible cross-tokenizer Bengali comparison are useful
to anyone building Bengali or Indic language technology.

\begin{thebibliography}{9}
\small
\bibitem{uax29} Unicode Consortium. \emph{Unicode Standard Annex \#29: Text Segmentation.} \url{https://unicode.org/reports/tr29/}
\bibitem{uax15} Unicode Consortium. \emph{Unicode Standard Annex \#15: Unicode Normalization Forms.} \url{https://unicode.org/reports/tr15/}
\bibitem{sarvam} Sarvam AI. \emph{Sarvam-1.} \url{https://www.sarvam.ai/blogs/sarvam-1}
\bibitem{indicsuper} \emph{IndicSuperTokenizer: An Optimized Tokenizer for Indic Multilingual LLMs.} arXiv:2511.03237.
\bibitem{indicgenbench} \emph{IndicGenBench.} arXiv:2404.16816.
\bibitem{bngraphemizer} \emph{Grapheme-level tokenization for Bengali handwriting recognition (BnGraphemizer).} PLATTER/CHIPS line of work.
\bibitem{brahmic} \emph{BrahmicTokenizer-131K.} arXiv:2605.29379.
\end{thebibliography}

\end{document}
